\documentclass{article}
\usepackage{graphicx}
\usepackage{amsmath}
\usepackage{amssymb}
\usepackage{amsfonts}
\usepackage{hyperref}
\usepackage{exercises}
\usepackage{indentfirst}
\usepackage{circledsteps}
\usepackage{amsthm}
\usepackage{csquotes}
\usepackage{yhmath}
\usepackage{comment}
\usepackage{placeins}
\usepackage{geometry}
\usepackage{hyperref}
\usepackage{tikz}
\newcommand*{\oldneg}{\mathord{\sim}}
\newcommand{\Z}{\mathbb{Z}}
\newcommand{\N}{\mathbb{N}}
\let\oldemptyset\emptyset

\theoremstyle{definition}
\newtheorem{definition}{Definition}
\newtheorem{theorem}{Theorem}[section]
\newtheorem{corollary}{Corollary}[theorem]
\newtheorem{lemma}[theorem]{Lemma}

\setlength{\parindent}{.5in}
\geometry{margin=0.5in}
\setlength{\parskip}{12pt}
\title{Chapter 8 Notes}
\begin{document}
\date{}
\maketitle

\section*{Section 1}

\subsection*{Notes}
Here is why, for any non-empty class of ordinals $B$, the smallest ordinal $\alpha\in B$ is an initial element of $B$. 
If $\alpha=\emptyset$, then $\alpha\cap B=\emptyset$. 
If $\alpha\ne\emptyset$, then all members of $\alpha$ are ordinals less than $\alpha$. But, by assumption, there are no ordinals in $B$ less than 
$\alpha$, meaning $\alpha\cap B=\emptyset$.

\subsection*{Exercise 1.1}
(1) We first show that $R_\Omega$ is well-founded. Consider any non-empty subclass $A$ of $R_\Omega$. Then let $\alpha$ be the 
lowest rank that any of the elements in $A$ have. Now consider any element $x\in A$ which has rank $\alpha$. 
Suppose that there was some element $y$ of $A$ in $x$. Then $y$ would have rank $\ge\alpha$. But this is not possible, as, by 
$P_6$ of Chapter 7 (and the fact that $x$ is a set), $y$ would have a rank $<\alpha$, a contradiction. Thus 
all elements of $x$ are not in $y$, meaning $x\cap A=\emptyset$. Thus every element $x\in A$ with rank $\alpha$ is an 
initial element of $A$.

(2) The fact that every $R_\alpha$ (where $\alpha>0$) is well-founded follows because every $R_\alpha$ is a non-empty subclass of 
$R_\Omega$. 

If now it was the case that $R_\Omega$ was in fact equal to $V$, this exercise would prove that 
$V$ was well-founded.

\subsection*{Theorem 2.3}
Here is a more formal proof of Theorem 2.3 (a more formal form of the second paragraph of the proof).

Since $B$ is non-empty, it contains an element $e$. We now form the denumerable sequence 
$b_1,b_2,...,b_n,...$ as follows: $b_1=e$; if $b_n$ is an initial element of $B$, then 
$b_{n+1}=b_n$; if $b_n$ is not an initial element of $B$, then we take $b_{n+1}$ to be the least element (under the well-ordering of $A$) of 
$b_n$ that is in $B$.

As remarked before the proof of Theorem 2.3, $A$ satisfying the descending $\in$-chain condition means that 
$A$ contains no $\in$-cycles. If all elements of the above 
denumerable sequence are distinct, then clearly that sequence is an infinite descending $\in$-chain, contradicting the assumption that 
$A$ satisfies the descending $\in$-chain condition. 

Suppose, on the other hand, that two elements of the sequence were the same 
(say the elements $b_m$ and $b_n$, where $m<n$), and that $b_m$ was not an initial element of $B$. Then, for all $i$ where $m\le i<n$, 
it must be that $b_i\in b_{i+1}$. But that means we would have an $\in$-cycle $b_m,b_{m+1},...,b_n$, which cannot happen. 

Thus it cannot be that all elements of the sequence are distinct, or that two elements of the sequence were the same where one of them 
was not an initial element of $B$. Hence it must be that a pair of elements of the sequence are the same, where one of the elements of the 
pair is an initial element.

\end{document}
